% chapters/ch02.tex
% 第~
\ref{chap:2}~章 一致有界原理与\index{bituxiangdingli@闭图像定理 (closed graph theorem)}
闭图像定理
% 源: 原书 PDF 第 41--64 页 (印刷页 31--54)

\chapter{一致有界原理与闭图像定理}
\label{chap:2}

\section{Baire 纲定理}
\label{sec:2.1}

下面的经典结果在第~
\ref{chap:2}~章各定理的证明中起着本质作用.

\begin{keytheorem}[Baire]
\label{thm:2.1}
设 $X$ 为完备度量空间, $(X_n)_{n\ge1}$ 为 $X$ 中的一列闭子集.假设对每个 $n\ge1$
都有
\begin{equation*}
\Int X_n=\emptyset.
\end{equation*}
则
\begin{equation*}
\Int\Bigl(\bigcup_{n=1}^\infty X_n\Bigr)=\emptyset.
\end{equation*}
\end{keytheorem}

\begin{remark}
\label{rem:2.1}
Baire 纲定理常以下述形式使用.设 $X$ 为非空完备度量空间, $(X_n)_{n\ge1}$ 为一列
闭子集, 使得
\begin{equation*}
\bigcup_{n=1}^\infty X_n=X.
\end{equation*}
则存在某个 $n_0$, 使 $\Int X_{n_0}\ne\emptyset$.
\end{remark}

\begin{proof}
置 $O_n=X_n^c$, 则对每个 $n\ge1$, $O_n$ 是 $X$ 中的开稠集.我们要证明
$G=\bigcap_{n=1}^\infty O_n$ 在 $X$ 中稠密.设 $\omega$ 为 $X$ 中的非空开集;
我们将证明 $\omega\cap G\ne\emptyset$.

如通常那样, 置
\begin{equation*}
B(x,r)=\{y\in X;\ d(y,x)<r\}.
\end{equation*}
任取 $x_0\in\omega$ 与 $r_0>0$, 使得
\begin{equation*}
B(x_0,r_0)\subset\omega.
\end{equation*}
然后取 $x_1\in B(x_0,r_0)\cap O_1$ 与 $r_1>0$, 使得
\begin{equation*}
\left\{\begin{aligned}
	&B(x_1,r_1)\subset B(x_0,r_0)\cap O_1,\\
	&0<r_1<\frac{r_0}{2},
\end{aligned}\right.
\end{equation*}
这总是可以做到的, 因为 $O_1$ 是开稠集.用归纳法构造两列 $(x_n)$ 与 $(r_n)$, 使得
\begin{equation*}
\left\{\begin{aligned}
	&B(x_{n+1},r_{n+1})\subset B(x_n,r_n)\cap O_{n+1},\quad \forall n\ge0,\\
	&0<r_{n+1}<\frac{r_n}{2}.
\end{aligned}\right.
\end{equation*}
由此可知 $(x_n)$ 是 Cauchy 列; 设 $x_n\to\ell$.

由于对每个 $n\ge0$ 与每个 $p\ge0$ 都有 $x_{n+p}\in B(x_n,r_n)$, 令 $p\to\infty$
取极限即得
\begin{equation*}
\ell\in\overline{B(x_n,r_n)},\quad \forall n\ge0.
\end{equation*}
特别地, $\ell\in\omega\cap G$.
\end{proof}

\section{一致有界原理}
\label{sec:2.2}

\begin{notation}
设 $E$ 与 $F$ 为两个赋范向量空间.以 $\LL(E,F)$ 记从 $E$ 到 $F$ 的连续 (即有界)
线性算子构成的空间, 赋以范数
\begin{equation*}
\norm{T}_{\LL(E,F)}=\sup_{\substack{x\in E\\ \norm{x}\le1}}\norm{Tx}.
\end{equation*}
如通常那样, 用 $\LL(E)$ 代替 $\LL(E,E)$.
\end{notation}

\index{BanachSteinhausdingli@Banach--Steinhaus 定理 (Banach--Steinhaus theorem)}
\begin{keytheorem}[Banach--Steinhaus, 一致有界原理]
\label{thm:2.2}
设 $E$ 与 $F$ 为两个 Banach 空间, $(T_i)_{i\in I}$ 为一族 (不必可数) 从 $E$ 到 $F$ 的
连续线性算子.假设
\begin{equation}
\sup_{i\in I}\norm{T_ix}<\infty\quad \forall x\in E.
\label{eq:2.1}
\end{equation}
则
\begin{equation}
\sup_{i\in I}\norm{T_i}_{\LL(E,F)}<\infty.
\label{eq:2.2}
\end{equation}
换言之, 存在常数 $c$ 使得
\begin{equation*}
\norm{T_ix}\le c\norm{x}\quad \forall x\in E,\ \forall i\in I.
\end{equation*}
\end{keytheorem}

\begin{remark}
\label{rem:2.2}
定理  \ref{thm:2.2} 的结论颇为出人意料.由逐点估计竟能推出整体 (一致) 估计.
\end{remark}

\begin{proof}
对每个 $n\ge1$, 令
\begin{equation*}
X_n=\{x\in E;\ \forall i\in I,\ \norm{T_ix}\le n\},
\end{equation*}
则 $X_n$ 是闭的, 且由 \eqref{eq:2.1} 有
\begin{equation*}
\bigcup_{n=1}^\infty X_n=E.
\end{equation*}
由 Baire 纲定理, 存在某个 $n_0\ge1$ 使 $\Int(X_{n_0})\ne\emptyset$.取 $x_0\in E$ 与
$r>0$ 使 $B(x_0,r)\subset X_{n_0}$.我们有
\begin{equation*}
\norm{T_i(x_0+rz)}\le n_0\quad \forall i\in I,\ \forall z\in B(0,1).
\end{equation*}
由此得到
\begin{equation*}
r\norm{T_i}_{\LL(E,F)}\le n_0+\norm{T_ix_0},
\end{equation*}
从而 \eqref{eq:2.2} 成立.
\end{proof}

\begin{remark}
\label{rem:2.3}
回忆: 一般说来, 连续映射的逐点极限不必连续.线性性假设在定理  \ref{thm:2.2} 中起着
本质作用.然而要注意: 在定理  \ref{thm:2.2} 的框架下, 并不能推出
$\norm{T_n-T}_{\LL(E,F)}\to0$.
\end{remark}

下面是一致有界原理的几个直接推论.

\begin{corollary}
\label{cor:2.3}
设 $E$ 与 $F$ 为两个 Banach 空间, $(T_n)$ 为从 $E$ 到 $F$ 的一列连续线性算子, 使得
对每个 $x\in E$, $T_nx$ 当 $n\to\infty$ 时收敛, 其极限记为 $Tx$.则
\begin{enumerate}[label=(\alph*), leftmargin=2.0em, itemsep=0.15em]
	\item $\sup_n\norm{T_n}_{\LL(E,F)}<\infty$,
	\item $T\in\LL(E,F)$,
	\item $\norm{T}_{\LL(E,F)}\le\liminf_{n\to\infty}\norm{T_n}_{\LL(E,F)}$.
\end{enumerate}
\end{corollary}

\begin{proof}
(a) 直接由定理  \ref{thm:2.2} 得到, 因而存在常数 $c$ 使得
\begin{equation*}
\norm{T_nx}\le c\norm{x}\quad \forall n,\ \forall x\in E.
\end{equation*}
取极限得
\begin{equation*}
\norm{Tx}\le c\norm{x}\quad \forall x\in E.
\end{equation*}
$T$ 显然是线性的, 故得 (b).

最后, 我们有
\begin{equation*}
\norm{T_nx}\le\norm{T_n}_{\LL(E,F)}\norm{x}\quad \forall x\in E,
\end{equation*}
于是 (c) 直接可得.
\end{proof}

\begin{keycorollary}
\label{cor:2.4}
设 $G$ 为 Banach 空间, $B$ 为 $G$ 的子集.假设
\begin{equation}
\text{对每个 } f\in G^\star,\ \text{集合 } f(B)=\{\dual{f}{x};\ x\in B\}
\text{ 有界 (在 } \R \text{ 中)}.
\label{eq:2.3}
\end{equation}
则
\begin{equation}
B \text{ 有界}.
\label{eq:2.4}
\end{equation}
\end{keycorollary}

\begin{proof}
我们对 $E=G^\star$, $F=\R$, $I=B$ 应用定理  \ref{thm:2.2}.对每个 $b\in B$, 置
\begin{equation*}
T_b(f)=\dual{f}{b},\quad f\in E=G^\star,
\end{equation*}
于是由 \eqref{eq:2.3},
\begin{equation*}
\sup_{b\in B}\abs{T_b(f)}<\infty\quad \forall f\in E.
\end{equation*}
由定理  \ref{thm:2.2}, 存在常数 $c$ 使得
\begin{equation*}
\abs{\dual{f}{b}}\le c\norm{f}\quad \forall f\in G^\star,\ \forall b\in B.
\end{equation*}
因此 (利用推论  \ref{cor:1.4}) 得到
\begin{equation*}
\norm{b}\le c\quad \forall b\in B.
\end{equation*}
\end{proof}

\begin{remark}
\label{rem:2.4}
推论  \ref{cor:2.4} 说的是: 要证明集合 $B$ 有界, 只须通过有界线性泛函去``观察''
$B$.这在有限维空间中是一种熟悉的做法, 那里线性泛函就是相对于某组基的分量.在
某种意义下, 推论  \ref{cor:2.4} 在无限维空间中取代了分量的用法.有时人们把推论
 \ref{cor:2.4} 的结论表述为: ``弱有界'' $\Longleftrightarrow$ ``强有界'' (见第~
\ref{chap:3}~章).
\end{remark}

下面是与推论  \ref{cor:2.4} 对偶的命题:

\begin{corollary}
\label{cor:2.5}
设 $G$ 为 Banach 空间, $B^\star$ 为 $G^\star$ 的子集.假设
\begin{equation}
\text{对每个 } x\in G,\ \text{集合 } \dual{B^\star}{x}=\{\dual{f}{x};\ f\in B^\star\}
\text{ 有界 (在 } \R \text{ 中)}.
\label{eq:2.5}
\end{equation}
则
\begin{equation}
B^\star \text{ 有界}.
\label{eq:2.6}
\end{equation}
\end{corollary}

\begin{proof}
对 $E=G$, $F=\R$, $I=B^\star$ 应用定理  \ref{thm:2.2}.对每个 $b\in B^\star$, 置
$T_b(x)=\dual{b}{x}$ ($x\in G=E$).我们得到: 存在常数 $c$ 使得
\begin{equation*}
\abs{\dual{b}{x}}\le c\norm{x}\quad \forall b\in B^\star,\ \forall x\in G.
\end{equation*}
由对偶范数的定义即得
\begin{equation*}
\norm{b}\le c\quad \forall b\in B^\star.
\end{equation*}
\end{proof}

\section{\index{kaiyingshedingli@开映射定理 (open mapping theorem)}
开映射定理与闭图像定理}
\label{sec:2.3}

下面是两个属于 Banach 的基本结果.

\begin{keytheorem}[开映射定理]
\label{thm:2.6}
设 $E$ 与 $F$ 为两个 Banach 空间, $T$ 为从 $E$ 到 $F$ 的连续线性算子且为满射
(= 映上).则存在常数 $c>0$ 使得
\begin{equation}
T(B_E(0,1))\supset B_F(0,c).
\label{eq:2.7}
\end{equation}
\end{keytheorem}

\begin{remark}
\label{rem:2.5}
性质 \eqref{eq:2.7} 蕴含: $T$ 把 $E$ 中任何开集的像都是 $F$ 中的开集 (这正是本
定理名称的由来!).事实上, 设 $U$ 在 $E$ 中开, 证明 $T(U)$ 开.任取
$y_0\in T(U)$, 则 $y_0=Tx_0$ 对某个 $x_0\in U$ 成立.取 $r>0$ 使
$B(x_0,r)\subset U$, 即 $x_0+B(0,r)\subset U$.于是
\begin{equation*}
y_0+T(B(0,r))\subset T(U).
\end{equation*}
利用 \eqref{eq:2.7} 得
\begin{equation*}
T(B(0,r))\supset B(0,rc),
\end{equation*}
从而
\begin{equation*}
B(y_0,rc)\subset T(U).
\end{equation*}
\end{remark}

定理  \ref{thm:2.6} 的几个重要推论如下.

\begin{keycorollary}
\index{youjienidingli@有界逆定理 (bounded inverse theorem)}
\label{cor:2.7}
设 $E$ 与 $F$ 为两个 Banach 空间, $T$ 为从 $E$ 到 $F$ 的连续线性算子且为双射,
即既单射 (= 一对一) 又满射.则 $T^{-1}$ 也连续 (从 $F$ 到 $E$).
\end{keycorollary}

\begin{proof}[推论 \ref{cor:2.7}]
性质 \eqref{eq:2.7} 与 $T$ 的单射性蕴含: 若 $x\in E$ 满足 $\norm{Tx}<c$, 则
$\norm{x}<1$.由齐次性得
\begin{equation*}
\norm{x}\le\frac{1}{c}\norm{Tx}\quad \forall x\in E,
\end{equation*}
因此 $T^{-1}$ 连续.
\end{proof}

\begin{corollary}
\label{cor:2.8}
设 $E$ 为赋以两个范数 $\norm{\cdot}_1$ 与 $\norm{\cdot}_2$ 的向量空间.假设 $E$ 对
这两个范数都是 Banach 空间, 且存在常数 $C\ge0$ 使得
\begin{equation*}
\norm{x}_2\le C\norm{x}_1\quad \forall x\in E.
\end{equation*}
则这两个范数等价, 即存在常数 $c>0$ 使得
\begin{equation*}
\norm{x}_1\le c\norm{x}_2\quad \forall x\in E.
\end{equation*}
\end{corollary}

\begin{proof}[推论 \ref{cor:2.8}]
对 $E=(E,\norm{\cdot}_1)$, $F=(E,\norm{\cdot}_2)$, $T=I$ 应用推论  \ref{cor:2.7}.
\end{proof}

\begin{proof}[定理 \ref{thm:2.6}]
我们把论证分为两步:

\emph{第一步.} 设 $T$ 为从 $E$ 映上 $F$ 的线性满射.则存在常数 $c>0$ 使得
\begin{equation}
T(B(0,1))\supset B(0,2c).
\label{eq:2.8}
\end{equation}

\emph{证明.} 置 $X_n=nT(B(0,1))$.因 $T$ 是满射, 有 $\bigcup_{n=1}^\infty X_n=F$;
由 Baire 纲定理, 存在某个 $n_0$ 使 $\Int(X_{n_0})\ne\emptyset$.由此
\begin{equation*}
\Int\bigl[T(B(0,1))\bigr]\ne\emptyset.
\end{equation*}
取 $c>0$ 与 $y_0\in F$ 使得
\begin{equation}
B(y_0,4c)\subset T(B(0,1)).
\label{eq:2.9}
\end{equation}
特别地, $y_0\in T(B(0,1))$; 又由对称性,
\begin{equation}
-y_0\in T(B(0,1)).
\label{eq:2.10}
\end{equation}
把 \eqref{eq:2.9} 与 \eqref{eq:2.10} 相加得
\begin{equation*}
B(0,4c)\subset T(B(0,1))+T(B(0,1)).
\end{equation*}
另一方面, 因 $T(B(0,1))$ 是凸的, 有
\begin{equation*}
T(B(0,1))+T(B(0,1))=2T(B(0,1)),
\end{equation*}
于是 \eqref{eq:2.8} 成立.

\emph{第二步.} 设 $T$ 为从 $E$ 到 $F$ 的连续线性算子且满足 \eqref{eq:2.8}.则
\begin{equation}
T(B(0,1))\supset B(0,c).
\label{eq:2.11}
\end{equation}

\emph{证明.} 任取 $y\in F$ 且 $\norm{y}<c$.目标是找 $x\in E$ 使得
$\norm{x}<1$ 且 $Tx=y$.由 \eqref{eq:2.8} 知
\begin{equation}
\forall\varepsilon>0\quad\exists z\in E\ \text{使}\ \norm{z}<\frac12\
\text{且}\ \norm{y-Tz}<\varepsilon.
\label{eq:2.12}
\end{equation}
取 $\varepsilon=c/2$, 找到 $z_1\in E$ 使得
\begin{equation*}
\norm{z_1}<\frac12\quad\text{且}\quad \norm{y-Tz_1}<\frac{c}{2}.
\end{equation*}
对 $y-Tz_1$ (代替 $y$) 施以同样的构造, 取 $\varepsilon=c/4$, 找到 $z_2\in E$ 使得
\begin{equation*}
\norm{z_2}<\frac14\quad\text{且}\quad \norm{(y-Tz_1)-Tz_2}<\frac{c}{4}.
\end{equation*}
如此继续, 用归纳法得到一列 $(z_n)$ 使得
\begin{equation*}
\norm{z_n}<\frac{1}{2^n}\quad\text{且}\quad
\norm{y-T(z_1+z_2+\cdots+z_n)}<\frac{c}{2^n}\quad \forall n.
\end{equation*}
由此可知数列 $x_n=z_1+z_2+\cdots+z_n$ 是 Cauchy 列.设 $x_n\to x$, 显然
$\norm{x}<1$ 且 $y=Tx$ (因 $T$ 连续).
\end{proof}

\begin{keytheorem}[闭图像定理]
\label{thm:2.9}
设 $E$ 与 $F$ 为两个 Banach 空间.设 $T$ 为从 $E$ 到 $F$ 的线性算子.假设 $T$ 的
图像 $G(T)$ 在 $E\times F$ 中闭.则 $T$ 连续.
\end{keytheorem}

\begin{remark}
\label{rem:2.6}
逆命题显然成立, 因为任何连续映射 (线性或非线性) 的图像都是闭的.
\end{remark}

\begin{proof}[定理 \ref{thm:2.9}]
在 $E$ 上考虑两个范数
\begin{equation*}
\norm{x}_1=\norm{x}_E+\norm{Tx}_F\quad\text{与}\quad \norm{x}_2=\norm{x}_E
\end{equation*}
(范数 $\norm{\cdot}_1$ 称为\emph{图像范数}).

利用 $G(T)$ 闭这一假设, 容易验证 $E$ 对范数 $\norm{\cdot}_1$ 是 Banach 空间.
另一方面, $E$ 对范数 $\norm{\cdot}_2$ 也是 Banach 空间, 且
$\norm{\cdot}_2\le\norm{\cdot}_1$.由推论  \ref{cor:2.8}, 这两个范数等价, 因而存在
常数 $c>0$ 使 $\norm{x}_1\le c\norm{x}_2$.于是 $\norm{Tx}_F\le c\norm{x}_E$.
\end{proof}

\starsection{补子空间. 线性算子的右可逆性与左可逆性}
\label{sec:2.4}

我们先给出 Banach 空间中闭子空间的若干几何性质, 它们可由开映射定理推出.

\begin{startheorem}
\label{thm:2.10}
设 $E$ 为 Banach 空间.假设 $G$ 与 $L$ 为两个闭线性子空间, 且 $G+L$ 闭.则存在
常数 $C\ge0$ 使得
\begin{equation}
\left\{\begin{aligned}
	&\text{每个 } z\in G+L \text{ 都可分解为}\\
	&z=x+y,\quad \text{其中 } x\in G,\ y\in L,\
	\norm{x}\le C\norm{z},\ \norm{y}\le C\norm{z}.
\end{aligned}\right.
\label{eq:2.13}
\end{equation}
\end{startheorem}

\begin{proof}
考虑乘积空间 $G\times L$, 赋以范数
\begin{equation*}
\norm{[x,y]}=\norm{x}+\norm{y},
\end{equation*}
以及赋以 $E$ 的范数的空间 $G+L$.

映射 $T:G\times L\to G+L$, $T[x,y]=x+y$, 是连续, 线性且满射的.由开映射定理,
存在常数 $c>0$, 使得每个满足 $\norm{z}<c$ 的 $z\in G+L$ 都可写成 $z=x+y$, 其中
$x\in G$, $y\in L$, 且 $\norm{x}+\norm{y}<1$.由齐次性, 每个 $z\in G+L$ 都可写成
\begin{equation*}
z=x+y\quad \text{其中 } x\in G,\ y\in L,\ \text{且}\
\norm{x}+\norm{y}\le(1/c)\norm{z}.
\end{equation*}
\end{proof}

\begin{starcorollary}
\label{cor:2.11}
在与定理  \ref{thm:2.10} 相同的假设下, 存在常数 $C$ 使得
\begin{equation}
\dist(x,G\cap L)\le C\{\dist(x,G)+\dist(x,L)\}\quad \forall x\in E.
\label{eq:2.14}
\end{equation}
\end{starcorollary}

\begin{proof}
给定 $x\in E$ 与 $\varepsilon>0$, 存在 $a\in G$ 与 $b\in L$ 使得
\begin{equation*}
\norm{x-a}\le\dist(x,G)+\varepsilon,\qquad
\norm{x-b}\le\dist(x,L)+\varepsilon.
\end{equation*}
把性质 \eqref{eq:2.13} 用于 $z=a-b$, 知存在 $a'\in G$ 与 $b'\in L$ 使得
\begin{equation*}
a-b=a'+b',\qquad \norm{a'}\le C\norm{a-b},\qquad \norm{b'}\le C\norm{a-b}.
\end{equation*}
于是 $a-a'\in G\cap L$, 且
\begin{align*}
\dist(x,G\cap L)&\le\norm{x-(a-a')}\le\norm{x-a}+\norm{a'}\\
&\le\norm{x-a}+C\norm{a-b}\le\norm{x-a}+C\bigl(\norm{x-a}+\norm{x-b}\bigr)\\
&\le(1+C)\dist(x,G)+\dist(x,L)+(1+2C)\varepsilon.
\end{align*}
最后令 $\varepsilon\to0$ 即得 \eqref{eq:2.14}.
\end{proof}

\begin{remark}
\label{rem:2.7}
推论  \ref{cor:2.11} 的逆命题也成立: 若 $G$ 与 $L$ 是两个闭线性子空间且
\eqref{eq:2.14} 成立, 则 $G+L$ 闭 (见习题  \ref{exer:2.16}).
\end{remark}

\begin{definition}
设 $G\subset E$ 为 Banach 空间 $E$ 的闭子空间.称子空间 $L\subset E$ 为 $G$ 的
\emph{拓扑\index{bukongjian@补空间 (complementary subspaces)}
补空间}, 或简称为 $G$ 的\emph{补空间}, 如果
\begin{enumerate}[label=(\roman*), leftmargin=2.4em, itemsep=0.15em]
	\item $L$ 闭,
	\item $G\cap L=\{0\}$ 且 $G+L=E$.
\end{enumerate}
此时也说 $G$ 与 $L$ 是 $E$ 的\emph{互补子空间}.若如此, 则每个 $z\in E$ 都可唯一地
写成 $z=x+y$, 其中 $x\in G$, $y\in L$.

由定理  \ref{thm:2.10} 可知, 投影算子 $z\mapsto x$ 与 $z\mapsto y$ 都是连续线性算子.
(这一性质也可以作为互补子空间的定义.)
\end{definition}

\subsection*{例}

\begin{enumerate}[label=\arabic*., leftmargin=2.2em, itemsep=0.5em]
	\item \emph{每个有限维子空间} $G$ 都有补空间.事实上, 设 $e_1,e_2,\dots,e_n$ 为
	      $G$ 的一组基.每个 $x\in G$ 可写成 $x=\sum_{i=1}^n x_ie_i$.置
	      $\varphi_i(x)=x_i$.利用 Hahn--Banach 定理 (解析形式) ——更确切地说, 推论
	       \ref{cor:1.2}——每个 $\varphi_i$ 都可延拓为定义在 $E$ 上的连续线性泛函
	      $\widetilde{\varphi}_i$.容易验证
	      $L=\bigcap_{i=1}^n(\widetilde{\varphi}_i)^{-1}(0)$ 是 $G$ 的补空间.
	\item \emph{每个有限余维的闭子空间} $G$ 都有补空间.只须取任一有限维空间 $L$,
	      使得 $G\cap L=\{0\}$ 且 $G+L=E$ ($L$ 因有限维而闭).
\end{enumerate}

下面是这类情形的一个典型例子.设 $N\subset E^\star$ 为 $p$ 维子空间.则
\begin{equation*}
G=\{x\in E;\ \dual{f}{x}=0\ \forall f\in N\}=N^\perp
\end{equation*}
是闭的, 且余维为 $p$.事实上, 设 $f_1,f_2,\dots,f_p$ 为 $N$ 的一组基.则存在
$e_1,e_2,\dots,e_p\in E$ 使得
\begin{equation*}
\dual{f_i}{e_j}=\delta_{ij}\quad \forall i,j=1,2,\dots,p.
\end{equation*}
[考虑由
\begin{equation}
\Phi(x)=\bigl(\dual{f_1}{x},\dual{f_2}{x},\dots,\dual{f_p}{x}\bigr)
\label{eq:2.15}
\end{equation}
定义的映射 $\Phi:E\to\R^p$, 注意 $\Phi$ 是满射; 否则由 Hahn--Banach 定理 (第二
几何形式), 会存在某个 $\alpha=(\alpha_1,\alpha_2,\dots,\alpha_p)\ne0$, 使得
\begin{equation*}
\alpha\cdot\Phi(x)=\Bigl\langle\sum_{i=1}^p\alpha_if_i,\ x\Bigr\rangle=0
\quad \forall x\in E,
\end{equation*}
这是荒谬的.]

容易验证向量组 $(e_i)_{1\le i\le p}$ 线性无关, 且由这些 $e_i$ 张成的空间是 $G$ 的
补空间.$N^\perp$ 的余维等于 $N$ 的维数这一事实的另一证明见第~
\ref{chap:11}~章 (命题~
\ref{prop:11.11}).

\begin{enumerate}[label=\arabic*., leftmargin=2.2em, itemsep=0.5em, start=3]
	\item 在 Hilbert 空间中, 每个闭子空间都有补空间 (见第~\ref{sec:5.2}~节).
\end{enumerate}

\begin{remark}
\label{rem:2.8}
重要的是要知道: 某些闭子空间 (即使在自反 Banach 空间中) \emph{没有}补空间.
事实上, J.~Lindenstrauss 与 L.~Tzafriri \bref{LindenstraussTzafriri}{1} 的一个著名
结果表明: 在每个不同构于 Hilbert 空间的 Banach 空间中, 都存在\emph{没有}补空间的
闭子空间.
\end{remark}

\begin{definition}
设 $T\in\LL(E,F)$.称算子 $S\in\LL(F,E)$ 为 $T$ 的\emph{右逆}, 如果
$T\circ S=I_F$.称算子 $S\in\LL(F,E)$ 为 $T$ 的\emph{左逆}, 如果
$S\circ T=I_E$.
\end{definition}

下面的结果给出这类逆存在的充分必要条件.

\begin{startheorem}
\label{thm:2.12}
假设 $T\in\LL(E,F)$ 是满射.则下述性质等价:
\begin{enumerate}[label=(\roman*), leftmargin=2.4em, itemsep=0.15em]
	\item $T$ 有右逆.
	\item $N(T)=T^{-1}(0)$ 在 $E$ 中有补空间.
\end{enumerate}
\end{startheorem}

\begin{proof}
(i) $\Rightarrow$ (ii). 设 $S$ 为 $T$ 的右逆.容易看出 (请验证) $R(S)=S(F)$ 是
$N(T)$ 在 $E$ 中的补空间.

(ii) $\Rightarrow$ (i). 设 $L$ 为 $N(T)$ 的补空间.设 $P$ 为从 $E$ 到 $L$ 的
(连续) 投影算子.给定 $f\in F$, 以 $x$ 记方程 $Tx=f$ 的任一解.置 $Sf=Px$.
注意 $S$ 与 $x$ 的选取无关.容易验证 $S\in\LL(F,E)$ 且 $T\circ S=I_F$.
\end{proof}

\begin{remark}
\label{rem:2.9}
由注记 \ref{rem:2.8} 与定理  \ref{thm:2.12}, 很容易构造没有右逆的满射算子 $T$.
事实上, 设 $G\subset E$ 为没有补空间的闭子空间, 令 $F=E/G$, 并令 $T$ 为从 $E$ 到
$F$ 的典范投影 (商空间的定义与性质见第~\ref{sec:11.2}~节).
\end{remark}

\begin{startheorem}
\label{thm:2.13}
假设 $T\in\LL(E,F)$ 是单射.则下述性质等价:
\begin{enumerate}[label=(\roman*), leftmargin=2.4em, itemsep=0.15em]
	\item $T$ 有左逆.
	\item $R(T)=T(E)$ 闭, 且在 $F$ 中有补空间.
\end{enumerate}
\end{startheorem}

\begin{proof}
(i) $\Rightarrow$ (ii). 容易验证 $R(T)$ 闭, 且 $N(S)$ 是 $R(T)$ 的补空间
[写 $f=TSf+(f-TSf)$].

(ii) $\Rightarrow$ (i). 设 $P$ 为从 $F$ 到 $R(T)$ 的连续投影算子.设 $f\in F$;
因 $Pf\in R(T)$, 存在唯一的 $x\in E$ 使 $Tx=Pf$.置 $Sf=x$.显然
$S\circ T=I_E$; 此外, 由推论  \ref{cor:2.7}, $S$ 连续.
\end{proof}

\starsection{再论正交性}
\label{sec:2.5}

有一些简单的公式, 给出和空间或交空间的正交表达.

\begin{proposition}
\label{prop:2.14}
设 $G$ 与 $L$ 为 $E$ 中的两个闭子空间.则
\begin{equation}
G\cap L=(G^\perp+L^\perp)^\perp,
\label{eq:2.16}
\end{equation}
\begin{equation}
G^\perp\cap L^\perp=(G+L)^\perp.
\label{eq:2.17}
\end{equation}
\end{proposition}

\begin{proof}[\eqref{eq:2.16}]
显然 $G\cap L\subset(G^\perp+L^\perp)^\perp$; 事实上, 若 $x\in G\cap L$ 且
$f\in G^\perp+L^\perp$, 则 $\dual{f}{x}=0$.反之, 我们有
$G^\perp\subset G^\perp+L^\perp$, 从而 $(G^\perp+L^\perp)^\perp\subset
G^{\perp\perp}=G$ (注意: 若 $N_1\subset N_2$, 则 $N_2^\perp\subset N_1^\perp$);
同理 $(G^\perp+L^\perp)^\perp\subset L$.因此 $(G^\perp+L^\perp)^\perp\subset
G\cap L$.
\end{proof}

\begin{proof}[\eqref{eq:2.17}]
用与 \eqref{eq:2.16} 的证明相同的论证.
\end{proof}

\begin{corollary}
\label{cor:2.15}
设 $G$ 与 $L$ 为 $E$ 中的两个闭子空间.则
\begin{equation}
(G\cap L)^\perp\supset G^\perp+L^\perp,
\label{eq:2.18}
\end{equation}
\begin{equation}
(G^\perp\cap L^\perp)^\perp=\overline{G+L}.
\label{eq:2.19}
\end{equation}
\end{corollary}

\begin{proof}
利用命题  \ref{prop:1.9} 与  \ref{prop:2.14}.
\end{proof}

下面是一个更深刻的结果.

\begin{startheorem}
\label{thm:2.16}
设 $G$ 与 $L$ 为 Banach 空间 $E$ 中的两个闭子空间.则下述性质等价:
\begin{enumerate}[label=(\alph*), leftmargin=2.0em, itemsep=0.15em]
	\item $G+L$ 在 $E$ 中闭,
	\item $G^\perp+L^\perp$ 在 $E^\star$ 中闭,
	\item $G+L=(G^\perp\cap L^\perp)^\perp$,
	\item $G^\perp+L^\perp=(G\cap L)^\perp$.
\end{enumerate}
\end{startheorem}

\begin{proof}
(a) $\Longleftrightarrow$ (c) 由 \eqref{eq:2.19} 得出.(d) $\Rightarrow$ (b) 是
显然的.

剩下的是蕴含关系 (a) $\Rightarrow$ (d) 与 (b) $\Rightarrow$ (a).

(a) $\Rightarrow$ (d). 由 \eqref{eq:2.18}, 只须证明
$(G\cap L)^\perp\subset G^\perp+L^\perp$.给定 $f\in(G\cap L)^\perp$, 考虑如下
定义的泛函 $\varphi:G+L\to\R$.对每个 $x\in G+L$, 写 $x=a+b$, 其中 $a\in G$, $b\in L$.置
\begin{equation*}
\varphi(x)=\dual{f}{a}.
\end{equation*}
显然 $\varphi$ 与 $x$ 的分解无关, 且 $\varphi$ 线性.另一方面, 由定理
 \ref{thm:2.10}, 我们可以选取 $x$ 的分解使得 $\norm{a}\le C\norm{x}$, 于是
\begin{equation*}
\abs{\varphi(x)}\le C\norm{x}\quad \forall x\in G+L.
\end{equation*}
把 $\varphi$ 延拓为整个 $E$ 上的连续线性泛函 $\widetilde{\varphi}$ (见推论
 \ref{cor:1.2}).于是
\begin{equation*}
f=(f-\widetilde{\varphi})+\widetilde{\varphi}
\quad\text{其中}\quad f-\widetilde{\varphi}\in G^\perp,\quad
\widetilde{\varphi}\in L^\perp.
\end{equation*}

(b) $\Rightarrow$ (a). 由推论  \ref{cor:2.11} 知存在常数 $C$ 使得
\begin{equation}
\dist(f,G^\perp\cap L^\perp)\le C\{\dist(f,G^\perp)+\dist(f,L^\perp)\}
\quad \forall f\in E^\star.
\label{eq:2.20}
\end{equation}
另一方面, 我们有
\begin{equation}
\dist(f,G^\perp)=\sup_{\substack{x\in G\\ \norm{x}\le1}}\dual{f}{x}
\quad \forall f\in E^\star.
\label{eq:2.21}
\end{equation}
[对 $\varphi(x)=I_{B_E}(x)-\dual{f}{x}$ 与 $\psi(x)=I_G(x)$ 应用定理
 \ref{thm:1.12}, 其中 $B_E=\{x\in E;\ \norm{x}\le1\}$.]

同理,
\begin{equation}
\dist(f,L^\perp)=\sup_{\substack{x\in L\\ \norm{x}\le1}}\dual{f}{x}
\quad \forall f\in E^\star,
\label{eq:2.22}
\end{equation}
并且 (由 \eqref{eq:2.17})
\begin{equation}
\dist(f,G^\perp\cap L^\perp)=\dist(f,(G+L)^\perp)
=\sup_{\substack{x\in G+L\\ \norm{x}\le1}}\dual{f}{x}\quad \forall f\in E^\star.
\label{eq:2.23}
\end{equation}
合并 \eqref{eq:2.20}, \eqref{eq:2.21}, \eqref{eq:2.22} 与 \eqref{eq:2.23} 得
\begin{equation}
\sup_{\substack{x\in G+L\\ \norm{x}\le1}}\dual{f}{x}
\le C\Bigl\{\sup_{\substack{x\in G\\ \norm{x}\le1}}\dual{f}{x}
+\sup_{\substack{x\in L\\ \norm{x}\le1}}\dual{f}{x}\Bigr\}
\quad \forall f\in E^\star.
\label{eq:2.24}
\end{equation}
由 \eqref{eq:2.24} 推出
\begin{equation}
\overline{B_G+B_L}\supset\frac{1}{C}B_{\overline{G+L}}.
\label{eq:2.25}
\end{equation}
事实上, 用反证法, 若存在 $x_0\in\overline{G+L}$ 满足 $\norm{x_0}\le1/C$ 且
$x_0\notin\overline{B_G+B_L}$, 则在 $E$ 中存在一张闭超平面把 $\{x_0\}$ 与
$\overline{B_G+B_L}$ 严格分离.于是存在 $f_0\in E^\star$ 与 $\alpha\in\R$ 使得
\begin{equation*}
\dual{f_0}{x}<\alpha<\dual{f_0}{x_0}\quad \forall x\in\overline{B_G+B_L}.
\end{equation*}
因此
\begin{equation*}
\sup_{\substack{x\in G\\ \norm{x}\le1}}\dual{f_0}{x}
+\sup_{\substack{x\in L\\ \norm{x}\le1}}\dual{f_0}{x}
\le\alpha<\dual{f_0}{x_0},
\end{equation*}
这与 \eqref{eq:2.24} 矛盾, 故 \eqref{eq:2.25} 得证.

最后, 考虑赋以范数
\begin{equation*}
\norm{[x,y]}=\max\{\norm{x},\norm{y}\}
\end{equation*}
的空间 $X=G\times L$, 以及赋以 $E$ 的范数的空间 $Y=\overline{G+L}$.由
$T([x,y])=x+y$ 定义的映射 $T:X\to Y$ 线性且连续.由 \eqref{eq:2.25} 知
\begin{equation*}
\overline{T(B_X)}\supset\frac{1}{C}B_Y.
\end{equation*}
利用定理  \ref{thm:2.6} (开映射定理) 证明中的第二步, 得到
\begin{equation*}
T(B_X)\supset\frac{1}{2C}B_Y.
\end{equation*}
由此可知 $T$ 从 $X$ 映上 $Y$, 即 $G+L=\overline{G+L}$.
\end{proof}

\section{\index{wujiexianxingsuanzi@无界线性算子 (unbounded linear operator)}
无界线性算子引论. \index{bansuisuanzi@伴随算子 (adjoint operator)}
伴随算子的定义}
\label{sec:2.6}

\begin{definition}
设 $E$ 与 $F$ 为两个 Banach 空间.从 $E$ 到 $F$ 的\emph{无界线性算子}是指定义在
线性子空间 $D(A)\subset E$ 上, 取值于 $F$ 的线性映射 $A:D(A)\subset E\to F$.集合
$D(A)$ 称为 $A$ 的\emph{定义域}.

若 $D(A)=E$ 且存在常数 $c\ge0$ 使
\begin{equation*}
\norm{Au}\le c\norm{u}\quad \forall u\in E,
\end{equation*}
则称 $A$ 是\emph{有界的} (或\emph{连续的}).

有界算子的范数定义为
\begin{equation*}
\norm{A}_{\LL(E,F)}=\operatorname*{Sup}_{u\ne0}\frac{\norm{Au}}{\norm{u}}.
\end{equation*}
\end{definition}

\begin{remark}
\label{rem:2.10}
当然可能出现这样的情况: 一个``无界线性算子''实际上是有界的.这个术语稍欠一致,
但通行且不会引起混淆.
\end{remark}

以下是一些重要的定义与进一步的记号:
\begin{equation*}
\boxed{\text{$A$ 的图像}=G(A)=\{[u,Au];\ u\in D(A)\}\subset E\times F,}
\end{equation*}
\begin{equation*}
\boxed{\text{$A$ 的值域}=R(A)=\{Au;\ u\in D(A)\}\subset F,}
\end{equation*}
\begin{equation*}
\boxed{\text{$A$ 的核}=N(A)=\{u\in D(A);\ Au=0\}\subset E.}
\end{equation*}
若 $G(A)$ 在 $E\times F$ 中闭, 则称映射 $A$ 是\emph{闭的}.

\begin{keyremark}
\label{rem:2.11}
要证明算子 $A$ 是闭的, 一般按如下步骤进行.取 $D(A)$ 中的序列 $(u_n)$, 使
$u_n\to u$ (在 $E$ 中) 且 $Au_n\to f$ (在 $F$ 中).然后验证两件事:
\begin{enumerate}[label=(\alph*), leftmargin=2.0em, itemsep=0.15em]
	\item $u\in D(A)$,
	\item $f=Au$.
\end{enumerate}
注意: 只考虑满足 $u_n\to0$ (在 $E$ 中) 且 $Au_n\to f$ (在 $F$ 中) 的序列
(并证明 $f=0$) 是\emph{不够}的.
\end{keyremark}

\begin{remark}
\label{rem:2.12}
若 $A$ 闭, 则 $N(A)$ 闭; 然而 $R(A)$ 不必闭.
\end{remark}

\begin{remark}
\label{rem:2.13}
实践中, 大多数无界算子是\emph{闭}的且是\emph{稠定}的, 即 $D(A)$ 在 $E$ 中稠密.
\end{remark}

\runin{伴随算子 $A^\star$ 的定义.}
设 $A:D(A)\subset E\to F$ 为稠定的无界线性算子.我们按下述方式引入无界算子
$A^\star:D(A^\star)\subset F^\star\to E^\star$.首先定义它的定义域:
\begin{equation*}
D(A^\star)=\{v\in F^\star;\ \exists c\ge0\ \text{使得}\
\abs{\dual{v}{Au}}\le c\norm{u}\ \forall u\in D(A)\}.
\end{equation*}
显然 $D(A^\star)$ 是 $F^\star$ 的线性子空间.现在定义 $A^\star v$.给定
$v\in D(A^\star)$, 考虑由
\begin{equation*}
g(u)=\dual{v}{Au}\quad \forall u\in D(A)
\end{equation*}
定义的映射 $g:D(A)\to\R$.我们有
\begin{equation*}
\abs{g(u)}\le c\norm{u}\quad \forall u\in D(A).
\end{equation*}
由 Hahn--Banach 定理 (解析形式; 见定理  \ref{thm:1.1}), 存在线性映射
$f:E\to\R$, 它延拓 $g$, 且满足
\begin{equation*}
\abs{f(u)}\le c\norm{u}\quad \forall u\in E.
\end{equation*}
于是 $f\in E^\star$.注意 $g$ 的延拓是唯一的, 因为 $D(A)$ 在 $E$ 中稠密.

置
\begin{equation*}
A^\star v=f.
\end{equation*}
这个无界线性算子 $A^\star:D(A^\star)\subset F^\star\to E^\star$ 称为 $A$ 的
\emph{伴随算子}.简言之, $A$ 与 $A^\star$ 之间的基本关系为
\begin{equation*}
\dual{v}{Au}_{F^\star,F}=\dual{A^\star v}{u}_{E^\star,E}
\quad \forall u\in D(A),\ \forall v\in D(A^\star).
\end{equation*}

\begin{remark}
\label{rem:2.14}
延拓 $g$ 不必借助 Hahn--Banach 定理.只须用经典的连续性延拓即可: 因为 $D(A)$
稠密, $g$ 在 $D(A)$ 上一致连续, 而 $\R$ 完备 (见例如 H.~L.~Royden
\bref{Royden}{1} (第~
\ref{chap:7}~章命题 11) 或 J.~Dugundji \bref{Dugundji}{1} (第 XIV 章
定理~
\ref{thm:5.2})).
\end{remark}

\begin{starremark}
\label{rem:2.15}
可能出现 $D(A^\star)$ 在 $F^\star$ 中不稠密的情形 (即使 $A$ 闭); 但这是一种相当
病态的情形 (见习题  \ref{exer:2.22}).下述断言总是成立: 若 $A$ 闭, 则
$D(A^\star)$ 对于第~
\ref{chap:3}~章定义的\wstar{}拓扑 $\sigma(F^\star,F)$ 在 $F^\star$ 中
稠密 (见问题 9).特别地, 若 $F$ 自反, 则 $D(A^\star)$ 对通常的 (范数) 拓扑在
$F^\star$ 中稠密 (见定理  \ref{thm:3.24}).
\end{starremark}

\begin{remark}
\label{rem:2.16}
若 $A$ 是有界算子, 则 $A^\star$ 也是有界算子 (从 $F^\star$ 到 $E^\star$), 并且
\begin{equation*}
\norm{A^\star}_{\LL(F^\star,E^\star)}=\norm{A}_{\LL(E,F)}.
\end{equation*}
事实上, 显然 $D(A^\star)=F^\star$.由基本关系,
\begin{equation*}
\abs{\dual{A^\star v}{u}}\le\norm{A}\norm{u}\norm{v}
\quad \forall u\in E,\ \forall v\in F^\star,
\end{equation*}
由此推出 $\norm{A^\star v}\le\norm{A}\norm{v}$, 从而 $\norm{A^\star}\le\norm{A}$.

我们还有
\begin{equation*}
\abs{\dual{v}{Au}}\le\norm{A^\star}\norm{u}\norm{v}
\quad \forall u\in E,\ \forall v\in F^\star,
\end{equation*}
由此 (由推论  \ref{cor:1.4}) 推出 $\norm{Au}\le\norm{A^\star}\norm{u}$, 从而
$\norm{A}\le\norm{A^\star}$.
\end{remark}

\begin{proposition}
\label{prop:2.17}
设 $A:D(A)\subset E\to F$ 为稠定的无界线性算子.则 $A^\star$ 闭, 即
$G(A^\star)$ 在 $F^\star\times E^\star$ 中闭.
\end{proposition}

\begin{proof}
设 $v_n\in D(A^\star)$ 使 $v_n\to v$ (在 $F^\star$ 中) 且 $A^\star v_n\to f$
(在 $E^\star$ 中).须验证 (a) $v\in D(A^\star)$ 与 (b) $A^\star v=f$.

我们有
\begin{equation*}
\dual{v_n}{Au}=\dual{A^\star v_n}{u}\quad \forall u\in D(A).
\end{equation*}
取极限得
\begin{equation*}
\dual{v}{Au}=\dual{f}{u}\quad \forall u\in D(A).
\end{equation*}
因此 $v\in D(A^\star)$ (因为
$\abs{\dual{v}{Au}}\le\norm{f}\norm{u}\ \forall u\in D(A)$), 且 $A^\star v=f$.
\end{proof}

\emph{$A$ 与 $A^\star$ 的图像由一个非常简单的正交关系相联系:} 考虑由
\begin{equation*}
I([v,f])=[-f,v]
\end{equation*}
定义的同构 $I:F^\star\times E^\star\to E^\star\times F^\star$.设
$A:D(A)\subset E\to F$ 为稠定的无界线性算子.则
\begin{equation*}
\boxed{I[G(A^\star)]=G(A)^\perp.}
\end{equation*}
事实上, 设 $[v,f]\in F^\star\times E^\star$, 则
\begin{align*}
[v,f]\in G(A^\star)
&\Longleftrightarrow \dual{f}{u}=\dual{v}{Au}\quad \forall u\in D(A)\\
&\Longleftrightarrow -\dual{f}{u}+\dual{v}{Au}=0\quad \forall u\in D(A)\\
&\Longleftrightarrow [-f,v]\in G(A)^\perp.
\end{align*}

下面是值域与核之间的若干标准正交关系:

\begin{corollary}
\label{cor:2.18}
设 $A:D(A)\subset E\to F$ 为稠定且闭的无界线性算子.则
\begin{align*}
\text{(i)}\qquad &N(A)=R(A^\star)^\perp,\\
\text{(ii)}\qquad &N(A^\star)=R(A)^\perp,\\
\text{(iii)}\qquad &N(A)^\perp\supset\overline{R(A^\star)},\\
\text{(iv)}\qquad &N(A^\star)^\perp=\overline{R(A)}.
\end{align*}
\end{corollary}

\begin{proof}
注意 (iii) 与 (iv) 直接由 (i), (ii) 结合命题  \ref{prop:1.9} 得到.(i) 与 (ii)
有简单直接的证明 (见习题  \ref{exer:2.18}).不过, 用下述办法把这些事实与命题
 \ref{prop:2.14} 联系起来是有启发性的.考虑空间 $X=E\times F$, 从而
$X^\star=E^\star\times F^\star$, 以及 $X$ 的子空间
\begin{equation*}
G=G(A)\quad\text{与}\quad L=E\times\{0\}.
\end{equation*}
很容易验证
\begin{equation}
N(A)\times\{0\}=G\cap L,
\label{eq:2.26}
\end{equation}
\begin{equation}
E\times R(A)=G+L,
\label{eq:2.27}
\end{equation}
\begin{equation}
\{0\}\times N(A^\star)=G^\perp\cap L^\perp,
\label{eq:2.28}
\end{equation}
\begin{equation}
R(A^\star)\times F^\star=G^\perp+L^\perp.
\label{eq:2.29}
\end{equation}

\emph{(i) 的证明.} 由 \eqref{eq:2.29} 有
\begin{equation*}
R(A^\star)^\perp\times\{0\}=(G^\perp+L^\perp)^\perp=G\cap L
\quad\text{(由 \eqref{eq:2.16})}
=N(A)\times\{0\}\quad\text{(由 \eqref{eq:2.26})}.
\end{equation*}

\emph{(ii) 的证明.} 由 \eqref{eq:2.27} 有
\begin{equation*}
\{0\}\times R(A)^\perp=(G+L)^\perp=G^\perp\cap L^\perp
\quad\text{(由 \eqref{eq:2.17})}
=\{0\}\times N(A^\star)\quad\text{(由 \eqref{eq:2.28})}.
\end{equation*}
\end{proof}

\begin{remark}
\label{rem:2.17}
即使 $A$ 是有界线性算子, 也可能出现 $N(A)^\perp\ne R(A^\star)$ (见习题
 \ref{exer:2.23}).不过, 下述断言总是成立: $N(A)^\perp$ 是 $R(A^\star)$ 关于
\wstar{}拓扑 $\sigma(E^\star,E)$ 的闭包 (见问题 9).特别地, 若 $E$ 自反, 则
$N(A)^\perp=R(A^\star)$.
\end{remark}

\starsection{闭值域算子的刻画. 满射算子的刻画}
\label{sec:2.7}

关于闭值域算子的主要结果如下.

\begin{startheorem}
\label{thm:2.19}
设 $A:D(A)\subset E\to F$ 为稠定且闭的无界线性算子.则下述性质等价:
\begin{enumerate}[label=(\roman*), leftmargin=2.4em, itemsep=0.15em]
	\item $R(A)$ 闭,
	\item $R(A^\star)$ 闭,
	\item $R(A)=N(A^\star)^\perp$,
	\item $R(A^\star)=N(A)^\perp$.
\end{enumerate}
\end{startheorem}

\begin{proof}
沿用推论  \ref{cor:2.18} 证明中的记号, 有
\begin{align*}
\text{(i)}&\Longleftrightarrow G+L\ \text{在 } X \text{ 中闭 (见 \eqref{eq:2.27})},\\
\text{(ii)}&\Longleftrightarrow G^\perp+L^\perp\ \text{在 } X^\star \text{ 中闭
(见 \eqref{eq:2.29})},\\
\text{(iii)}&\Longleftrightarrow G+L=(G^\perp\cap L^\perp)^\perp
\ \text{(见 \eqref{eq:2.27} 与 \eqref{eq:2.28})},\\
\text{(iv)}&\Longleftrightarrow (G\cap L)^\perp=G^\perp+L^\perp
\ \text{(见 \eqref{eq:2.26} 与 \eqref{eq:2.29})}.
\end{align*}
结论于是由定理  \ref{thm:2.16} 得到.
\end{proof}

\begin{remark}
\label{rem:2.18}
设 $A:D(A)\subset E\to F$ 为闭的无界线性算子.则 $R(A)$ 闭当且仅当存在常数 $C$
使得
\begin{equation*}
\dist(u,N(A))\le C\norm{Au}\quad \forall u\in D(A);
\end{equation*}
见习题  \ref{exer:2.14}.
\end{remark}

下一个结果给出满射算子的一个有用刻画.

\begin{startheorem}
\label{thm:2.20}
设 $A:D(A)\subset E\to F$ 为稠定且闭的无界线性算子.则下述性质等价:
\begin{enumerate}[label=(\alph*), leftmargin=2.0em, itemsep=0.15em]
	\item $A$ 是满射, 即 $R(A)=F$,
	\item 存在常数 $C$ 使得
	      \begin{equation*}
	      \norm{v}\le C\norm{A^\star v}\quad \forall v\in D(A^\star),
	      \end{equation*}
	\item $N(A^\star)=\{0\}$ 且 $R(A^\star)$ 闭.
\end{enumerate}
\end{startheorem}

\begin{remark}
\label{rem:2.19}
蕴含关系 (b) $\Rightarrow$ (a) 在实践中有时很有用, 可用来断定算子 $A$ 是满射.
做法如下.假设 $v$ 满足 $A^\star v=f$, 试证 $\norm{v}\le C\norm{f}$ (其中 $C$
与 $f$ 无关).这称为\emph{先验估计}方法.人们并不关心方程 $A^\star v=f$ 是否
有解; 而是先验地假定 $v$ 已给定, 再设法估计其范数.
\end{remark}

\begin{proof}
(a) $\Rightarrow$ (b). 置
\begin{equation*}
B^\star=\{v\in D(A^\star);\ \norm{A^\star v}\le1\}.
\end{equation*}
由齐次性, 只须证明 $B^\star$ 有界.为此——由推论  \ref{cor:2.5} (\index{yizhiyoujieyuanli@一致有界原理 (uniform boundedness principle)}
一致有界原理)
——只须证明: 对任意给定的 $f_0\in F$, 集合 $\dual{B^\star}{f_0}$ 有界 (在 $\R$
中).因 $A$ 是满射, 存在 $u_0\in D(A)$ 使 $Au_0=f_0$.对每个 $v\in B^\star$ 有
\begin{equation*}
\dual{v}{f_0}=\dual{v}{Au_0}=\dual{A^\star v}{u_0},
\end{equation*}
从而 $\abs{\dual{v}{f_0}}\le\norm{u_0}$.

(b) $\Rightarrow$ (c). 设 $f_n=A^\star v_n\to f$.对 $v_n-v_m$ 使用 (b) 可知
$(v_n)$ 是 Cauchy 列, 于是 $v_n\to v$.因 $A^\star$ 闭 (由命题  \ref{prop:2.17}),
得 $A^\star v=f$.

(c) $\Rightarrow$ (a). 因 $R(A^\star)$ 闭, 由定理  \ref{thm:2.19} 得
$R(A)=N(A^\star)^\perp=F$.
\end{proof}

还有一个``对偶''命题.

\begin{startheorem}
\label{thm:2.21}
设 $A:D(A)\subset F$ 为稠定且闭的无界线性算子.则下述性质等价:
\begin{enumerate}[label=(\alph*), leftmargin=2.0em, itemsep=0.15em]
	\item $A^\star$ 是满射, 即 $R(A^\star)=E^\star$,
	\item 存在常数 $C$ 使得
	      \begin{equation*}
	      \norm{u}\le C\norm{Au}\quad \forall u\in D(A),
	      \end{equation*}
	\item $N(A)=\{0\}$ 且 $R(A)$ 闭.
\end{enumerate}
\end{startheorem}

\begin{proof}
证明与定理  \ref{thm:2.20} 类似, 留作练习.
\end{proof}

\begin{remark}
\label{rem:2.20}
若假设 $\dim E<\infty$ 或 $\dim F<\infty$, 则下述各条等价:
\begin{align*}
&A \text{ 满射} \Longleftrightarrow A^\star \text{ 单射},\\
&A^\star \text{ 满射} \Longleftrightarrow A \text{ 单射},
\end{align*}
这确实是有限维空间中线性算子的经典结果.这些等价关系成立的原因是 $R(A)$ 与
$R(A^\star)$ 都是有限维的 (从而闭).

在一般情形下, 只有蕴含关系
\begin{align*}
&A \text{ 满射} \Rightarrow A^\star \text{ 单射},\\
&A^\star \text{ 满射} \Rightarrow A \text{ 单射}.
\end{align*}
逆命题不成立, 从下面这个简单的例子可以看出.设 $E=F=\ell^2$; 对每个
$x\in\ell^2$ 写 $x=(x_n)_{n\ge1}$, 并置 $Ax=\bigl(\frac1n x_n\bigr)_{n\ge1}$.
容易看出 $A$ 是有界算子且 $A^\star=A$; $A^\star$ (相应地 $A$) 是单射, 但 $A$
(相应地 $A^\star$) 不是满射; $R(A)$ (相应地 $R(A^\star)$) 稠密但不闭.
\end{remark}

\section*{第 2 章评注}
\addcontentsline{toc}{section}{第 2 章评注}
\markright{第 2 章评注}

\cmtnum{1.}
可以显式地写出一些没有补空间的简单闭子空间.例如 $c_0$ 是 $\ell^\infty$ 中
没有补空间的闭子空间; 见例如 C.~DeVito \bref{DeVito}{1} (记号 $c_0$ 与
$\ell^\infty$ 的说明见第~\ref{sec:11.3}~节).其他例子见 W.~Rudin \bref{Rudin}{1} ($L^1$
的一个子空间), G.~Köthe \bref{Kothe}{1} 以及 B.~Beauzamy \bref{Beauzamy}{1}
($\ell^p$ 的一个子空间, $p\ne2$).

\cmtnum{2.}
第~
\ref{chap:2}~章的大部分结果都可推广到 \emph{Fréchet 空间} (可度量化且完备的局部凸
空间).可能的推广方式很多; 见例如 H.~Schaefer \bref{Schaefer}{1}, J.~Horváth \bref{Horvath}{1}, R.~Edwards \bref{Edwards}{1}, F.~Treves
\bref{Treves}{1}, \bref{Treves}{3}, G.~Köthe \bref{Kothe}{1}.这些推广的动机
来自\emph{广义函数理论} (见 L.~Schwartz \bref{Schwartz}{1}), 其中许多重要的空间
\emph{不是} Banach 空间.关于在偏微分方程理论中的应用, 读者可参阅
L.~Hörmander \bref{Hormander}{1} 或 F.~Treves \bref{Treves}{1}, \bref{Treves}{2},
\bref{Treves}{3}.

\cmtnum{3.}
第~\ref{sec:2.5}~节的结果在 T.~Kato \bref{Kato}{1} 中有各种推广.

\section*{第 2 章习题}
\addcontentsline{toc}{section}{第 2 章习题}
\markright{第 2 章习题}

\begin{exercise}[凸函数的连续性.]
\label{exer:2.1}
设 $E$ 为 Banach 空间, $\varphi:E\to(-\infty,+\infty]$ 为凸的 l.s.c. 函数.
假设 $x_0\in\Int D(\varphi)$.

\begin{enumerate}[label=\arabic*., leftmargin=2.2em, itemsep=0.25em]
	\item 证明存在两个常数 $R>0$ 与 $M$, 使得
	      \begin{equation*}
	      \varphi(x)\le M\quad \forall x\in E\ \text{满足}\ \norm{x-x_0}\le R.
	      \end{equation*}
	      [提示: 取适当的 $\rho>0$, 考虑集合
	      \begin{equation*}
	      F_n=\{x\in E;\ \norm{x-x_0}\le\rho\ \text{且}\ \varphi(x)\le n\}.]
	      \end{equation*}
	\item 证明: $\forall r<R$, $\exists L\ge0$ 使得
	      \begin{equation*}
	      \abs{\varphi(x_1)-\varphi(x_2)}\le L\norm{x_1-x_2}
	      \quad \forall x_1,x_2\in E\ \text{满足}\
	      \norm{x_i-x_0}\le r,\ i=1,2.
	      \end{equation*}
	      更精确地说, 可以取 $L=\dfrac{2[M-\varphi(x_0)]}{R-r}$.
\end{enumerate}
\end{exercise}

\begin{exercise}
\label{exer:2.2}
设 $E$ 为向量空间, $p:E\to\R$ 为具有下述三条性质的函数:
\begin{enumerate}[label=(\roman*), leftmargin=2.4em, itemsep=0.15em]
	\item $p(x+y)\le p(x)+p(y)\quad \forall x,y\in E$,
	\item 对每个固定的 $x\in E$, 函数 $\lambda\mapsto p(\lambda x)$ 是从 $\R$ 到
	      $\R$ 的连续函数,
	\item 只要 $E$ 中的序列 $(y_n)$ 满足 $p(y_n)\to0$, 就有
	      $p(\lambda y_n)\to0$ 对每个 $\lambda\in\R$ 成立.
\end{enumerate}
假设 $(x_n)$ 为 $E$ 中的序列, 满足 $p(x_n)\to0$, 而 $(\alpha_n)$ 为 $\R$ 中的
有界序列.证明 $p(0)=0$ 且 $p(\alpha_nx_n)\to0$.

[提示: 给定 $\varepsilon>0$, 考虑集合
\begin{equation*}
F_n=\{\lambda\in\R;\ \abs{p(\lambda x_k)}\le\varepsilon,\ \forall k\ge n\}.]
\end{equation*}

并由此推出: 若 $(x_n)$ 为 $E$ 中的序列, 满足 $p(x_n-x)\to0$ 对某个 $x\in E$
成立, 而 $(\alpha_n)$ 为 $\R$ 中的序列, 满足 $\alpha_n\to\alpha$, 则
$p(\alpha_nx_n)\to p(\alpha x)$.
\end{exercise}

\begin{exercise}
\label{exer:2.3}
设 $E$ 与 $F$ 为两个 Banach 空间, $(T_n)$ 为 $\LL(E,F)$ 中的序列.假设对每个
$x\in E$, $T_nx$ 当 $n\to\infty$ 时收敛, 其极限记为 $Tx$.证明: 若 $x_n\to x$
(在 $E$ 中), 则 $T_nx_n\to Tx$ (在 $F$ 中).
\end{exercise}

\begin{exercise}
\label{exer:2.4}
设 $E$ 与 $F$ 为两个 Banach 空间, $a:E\times F\to\R$ 为满足下述条件的双线性
形式:
\begin{enumerate}[label=(\roman*), leftmargin=2.4em, itemsep=0.15em]
	\item 对每个固定的 $x\in E$, 映射 $y\mapsto a(x,y)$ 连续;
	\item 对每个固定的 $y\in F$, 映射 $x\mapsto a(x,y)$ 连续.
\end{enumerate}
证明存在常数 $C\ge0$ 使得
\begin{equation*}
\abs{a(x,y)}\le C\norm{x}\norm{y}\quad \forall x\in E,\ \forall y\in F.
\end{equation*}
[提示: 引入线性算子 $T:E\to F^\star$, 并借助推论  \ref{cor:2.5} 证明 $T$ 有界.]
\end{exercise}

\begin{exercise}
\label{exer:2.5}
设 $E$ 为 Banach 空间, $\varepsilon_n$ 为一列正数, 满足 $\lim\varepsilon_n=0$.
又设 $(f_n)$ 为 $E^\star$ 中的序列, 具有下述性质:
\begin{equation*}
\left\{\begin{aligned}
	&\exists r>0,\quad \forall x\in E\ \text{且}\ \norm{x}<r,\ \exists C(x)\in\R\
	\text{使得}\\
	&\dual{f_n}{x}\le\varepsilon_n\norm{f_n}+C(x)\quad \forall n.
\end{aligned}\right.
\end{equation*}
证明 $(f_n)$ 有界.

[提示: 引入 $g_n=f_n/(1+\varepsilon_n\norm{f_n})$.]
\end{exercise}

\begin{exercise}[局部有界的非线性单调算子.]
\label{exer:2.6}
设 $E$ 为 Banach 空间, $D(A)$ 为 $E$ 的任意子集.称 (非线性) 映射
$A:D(A)\subset E\to E^\star$ 为\emph{单调}的, 如果它满足
\begin{equation*}
\dual{Ax-Ay}{x-y}\ge0\quad \forall x,y\in D(A).
\end{equation*}

\begin{enumerate}[label=\arabic*., leftmargin=2.2em, itemsep=0.25em]
	\item 设 $x_0\in\Int D(A)$.证明存在两个常数 $R>0$ 与 $C$, 使得
	      \begin{equation*}
	      \norm{Ax}\le C\quad \forall x\in D(A)\ \text{满足}\ \norm{x-x_0}<R.
	      \end{equation*}
	      [提示: 用反证法, 构造 $D(A)$ 中的序列 $(x_n)$, 使得 $x_n\to x_0$ 且
	      $\norm{Ax_n}\to\infty$.取 $r>0$ 使 $B(x_0,r)\subset D(A)$.利用 $A$
	      在 $x_n$ 处与在 $(x_0+x)$ 处 (其中 $\norm{x}<r$) 的单调性.应用习题
	       \ref{exer:2.5}.]
	\item 对点 $x_0\in\Int[\operatorname{conv}D(A)]$ 证明同样的结论.
	\item 把第 1 问的结论推广到 $A$ 为多值的情形, 即对每个 $x\in D(A)$, $Ax$
	      是 $E^\star$ 的非空子集; 单调性定义如下:
	      \begin{equation*}
	      \dual{f-g}{x-y}\ge0\quad \forall x,y\in D(A),\ \forall f\in Ax,\
	      \forall g\in Ay.
	      \end{equation*}
\end{enumerate}
\end{exercise}

\begin{exercise}
\label{exer:2.7}
设 $\alpha=(\alpha_n)$ 为给定的实数列, $1\le p\le\infty$.假设对 $\ell^p$ 中
的每个元素 $x=(x_n)$ 都有 $\sum\abs{\alpha_n}\abs{x_n}<\infty$ ($\ell^p$ 空间的
定义见第~\ref{sec:11.3}~节).

证明 $\alpha\in\ell^{p'}$.
\end{exercise}

\begin{exercise}
\label{exer:2.8}
设 $E$ 为 Banach 空间, $T:E\to E^\star$ 为满足
\begin{equation*}
\dual{Tx}{x}\ge0\quad \forall x\in E
\end{equation*}
的线性算子.证明 $T$ 是有界算子.

[有两种方法: (i) 利用习题  \ref{exer:2.6}; 或 (ii) 应用闭图像定理.]
\end{exercise}

\begin{exercise}
\label{exer:2.9}
设 $E$ 为 Banach 空间, $T:E\to E^\star$ 为满足
\begin{equation*}
\dual{Tx}{y}=\dual{Ty}{x}\quad \forall x,y\in E
\end{equation*}
的线性算子.证明 $T$ 是有界算子.
\end{exercise}

\begin{exercise}
\label{exer:2.10}
设 $E$ 与 $F$ 为两个 Banach 空间, $T\in\LL(E,F)$ 为满射.

\begin{enumerate}[label=\arabic*., leftmargin=2.2em, itemsep=0.25em]
	\item 设 $M$ 为 $E$ 的任意子集.证明: $T(M)$ 在 $F$ 中闭当且仅当
	      $M+N(T)$ 在 $E$ 中闭.
	\item 由此推出: 若 $M$ 为 $E$ 中的闭向量空间且 $\dim N(T)<\infty$, 则
	      $T(M)$ 闭.
\end{enumerate}
\end{exercise}

\begin{exercise}
\label{exer:2.11}
设 $E$ 为 Banach 空间, $F=\ell^1$, $T\in\LL(E,F)$ 为满射.证明存在
$S\in\LL(F,E)$ 使得 $T\circ S=I_F$, 即 $S$ 是 $T$ 的右逆.

[提示: 不要用定理  \ref{thm:2.12}; 试着利用 $\ell^1$ 的典范基显式地定义 $S$.]
\end{exercise}

\begin{exercise}
\label{exer:2.12}
设 $E$ 与 $F$ 为两个 Banach 空间, 其范数分别为 $\norm{\cdot}_E$ 与
$\norm{\cdot}_F$.设 $T\in\LL(E,F)$ 使得 $R(T)$ 闭且 $\dim N(T)<\infty$.以
$\abs{\cdot}$ 记 $E$ 上的另一个范数, 它比 $\norm{\cdot}_E$ 弱, 即
$\abs{x}\le M\norm{x}_E\ \forall x\in E$.

证明存在常数 $C$ 使得
\begin{equation*}
\norm{x}_E\le C\bigl(\norm{Tx}_F+\abs{x}\bigr)\quad \forall x\in E.
\end{equation*}
[提示: 用反证法.]
\end{exercise}

\begin{exercise}
\label{exer:2.13}
设 $E$ 与 $F$ 为两个 Banach 空间.证明集合
\begin{equation*}
\Omega=\{T\in\LL(E,F);\ T\ \text{有左逆}\}
\end{equation*}
在 $\LL(E,F)$ 中开.

[提示: 先证明集合
\begin{equation*}
\mathcal{O}=\{T\in\LL(E,F);\ T\ \text{是双射}\}
\end{equation*}
在 $\LL(E,F)$ 中开.]
\end{exercise}

\begin{exercise}
\label{exer:2.14}
设 $E$ 与 $F$ 为两个 Banach 空间.

\begin{enumerate}[label=\arabic*., leftmargin=2.2em, itemsep=0.25em]
	\item 设 $T\in\LL(E,F)$.证明 $R(T)$ 闭当且仅当存在常数 $C$ 使得
	      \begin{equation*}
	      \dist(x,N(T))\le C\norm{Tx}\quad \forall x\in E.
	      \end{equation*}
	      [提示: 利用商空间 $E/N(T)$; 见第~\ref{sec:11.2}~节.]
	\item 设 $A:D(A)\subset E\to F$ 为闭的无界算子.证明 $R(A)$ 闭当且仅当
	      存在常数 $C$ 使得
	      \begin{equation*}
	      \dist(u,N(A))\le C\norm{Au}\quad \forall u\in D(A).
	      \end{equation*}
	      [提示: 考虑算子 $T:E_0\to F$, 其中 $E_0=D(A)$ 赋以图像范数, 且
	      $T=A$.]
\end{enumerate}
\end{exercise}

\begin{exercise}
\label{exer:2.15}
设 $E_1$, $E_2$ 与 $F$ 为三个 Banach 空间.设 $T_1\in\LL(E_1,F)$,
$T_2\in\LL(E_2,F)$ 使得
\begin{equation*}
R(T_1)\cap R(T_2)=\{0\}\quad\text{且}\quad R(T_1)+R(T_2)=F.
\end{equation*}
证明 $R(T_1)$ 与 $R(T_2)$ 都闭.

[提示: 对由 $T(x_1,x_2)=T_1x_1+T_2x_2$ 定义的映射
$T:E_1\times E_2\to F$ 应用习题  \ref{exer:2.10}.]
\end{exercise}

\begin{exercise}
\label{exer:2.16}
设 $E$ 为 Banach 空间, $G$ 与 $L$ 为 $E$ 的两个闭子空间.假设存在常数 $C$
使得
\begin{equation*}
\dist(x,G\cap L)\le C\dist(x,L),\quad \forall x\in G.
\end{equation*}
证明 $G+L$ 闭.
\end{exercise}

\begin{exercise}
\label{exer:2.17}
设 $E=C([0,1])$, 赋以通常的范数.考虑由
\begin{equation*}
D(A)=C^1([0,1])\quad\text{与}\quad Au=u'=\frac{\dd u}{\dd t}
\end{equation*}
定义的算子 $A:D(A)\subset E\to E$.

\begin{enumerate}[label=\arabic*., leftmargin=2.2em, itemsep=0.25em]
	\item 验证 $\overline{D(A)}=E$.
	\item $A$ 是否闭?
	\item 考虑由
	      \begin{equation*}
	      D(B)=C^2([0,1])\quad\text{与}\quad Bu=u'=\frac{\dd u}{\dd t}
	      \end{equation*}
	      定义的算子 $B:D(B)\subset E\to E$.$B$ 是否闭?
\end{enumerate}
\end{exercise}

\begin{exercise}
\label{exer:2.18}
设 $E$ 与 $F$ 为两个 Banach 空间, $A:D(A)\subset E\to F$ 为稠定的无界算子.

\begin{enumerate}[label=\arabic*., leftmargin=2.2em, itemsep=0.25em]
	\item 证明 $N(A^\star)=R(A)^\perp$ 且 $N(A)\subset R(A^\star)^\perp$.
	\item 假设 $A$ 还是闭的, 证明 $N(A)=R(A^\star)^\perp$.
\end{enumerate}
[试着找出直接的论证, 不要依赖推论  \ref{cor:2.18} 的证明.第 2 问用反证法:
假设存在 $u\in R(A^\star)^\perp$ 使 $[u,0]\notin G(A)$, 并应用 Hahn--Banach 定理.]
\end{exercise}

\begin{exercise}
\label{exer:2.19}
设 $E$ 为 Banach 空间, $A:D(A)\subset E\to E^\star$ 为稠定的无界算子.

\begin{enumerate}[label=\arabic*., leftmargin=2.2em, itemsep=0.25em]
	\item 假设存在常数 $C$ 使得
	      \begin{equation}
	      \dual{Au}{u}\ge-C\norm{Au}^2\quad \forall u\in D(A).
	      \tag{1}
	      \end{equation}
	      证明 $N(A)\subset N(A^\star)$.
	\item 反之, 假设 $N(A)\subset N(A^\star)$.又假设 $A$ 闭且 $R(A)$ 闭.
	      证明存在常数 $C$ 使 (1) 成立.
\end{enumerate}
\end{exercise}

\begin{exercise}
\label{exer:2.20}
设 $E$ 与 $F$ 为两个 Banach 空间.设 $T\in\LL(E,F)$, $A:D(A)\subset E\to F$
为稠定且闭的无界算子.考虑由
\begin{equation*}
D(B)=D(A),\qquad B=A+T
\end{equation*}
定义的算子 $B:D(B)\subset E\to F$.

\begin{enumerate}[label=\arabic*., leftmargin=2.2em, itemsep=0.25em]
	\item 证明 $B$ 闭.
	\item 证明 $D(B^\star)=D(A^\star)$ 且 $B^\star=A^\star+T^\star$.
\end{enumerate}
\end{exercise}

\begin{exercise}
\label{exer:2.21}
设 $E$ 为无限维 Banach 空间.固定元素 $a\in E$, $a\ne0$, 以及不连续的线性
泛函 $f:E\to\R$ (这样的泛函存在; 见习题  \ref{exer:1.5}).考虑由
\begin{equation*}
D(A)=E,\qquad Ax=x-f(x)a
\end{equation*}
定义的算子 $A:E\to E$.

\begin{enumerate}[label=\arabic*., leftmargin=2.2em, itemsep=0.25em]
	\item 确定 $N(A)$ 与 $R(A)$.
	\item $A$ 是否闭?
	\item 确定 $A^\star$ (仔细地定义 $D(A^\star)$).
	\item 确定 $N(A^\star)$ 与 $R(A^\star)$.
	\item 比较 $N(A)$ 与 $R(A^\star)^\perp$, 以及 $N(A^\star)$ 与 $R(A)^\perp$.
	\item 与习题  \ref{exer:2.18} 的结果作比较.
\end{enumerate}
\end{exercise}

\begin{exercise}
\label{exer:2.22}
本习题的目的是构造一个无界算子 $A:D(A)\subset E\to E$, 它是稠定, 闭的, 且
满足 $\overline{D(A^\star)}\ne E^\star$.

设 $E=\ell^1$, 从而 $E^\star=\ell^\infty$.考虑由
\begin{equation*}
D(A)=\bigl\{u=(u_n)\in\ell^1;\ (nu_n)\in\ell^1\bigr\}
\quad\text{与}\quad Au=(nu_n)
\end{equation*}
定义的算子 $A:D(A)\subset E\to E$.

\begin{enumerate}[label=\arabic*., leftmargin=2.2em, itemsep=0.25em]
	\item 验证 $A$ 是稠定且闭的.
	\item 确定 $D(A^\star)$, $A^\star$ 与 $\overline{D(A^\star)}$.
\end{enumerate}
\end{exercise}

\begin{exercise}
\label{exer:2.23}
设 $E=\ell^1$, 从而 $E^\star=\ell^\infty$.考虑由
\begin{equation*}
Tu=\Bigl(\frac1n u_n\Bigr)_{n\ge1}
\quad\text{对每个 } u=(u_n)_{n\ge1}\in\ell^1
\end{equation*}
定义的算子 $T\in\LL(E,E)$.

确定 $N(T)$, $N(T)^\perp$, $T^\star$, $R(T^\star)$ 与
$\overline{R(T^\star)}$.与推论  \ref{cor:2.18} 作比较.
\end{exercise}

\begin{exercise}
\label{exer:2.24}
设 $E$, $F$ 与 $G$ 为三个 Banach 空间.设 $A:D(A)\subset E\to F$ 为稠定的
无界算子.设 $T\in\LL(F,G)$, 考虑由 $D(B)=D(A)$ 与 $B=T\circ A$ 定义的算子
$B:D(B)\subset E\to G$.

\begin{enumerate}[label=\arabic*., leftmargin=2.2em, itemsep=0.25em]
	\item 确定 $B^\star$.
	\item (用一个例子) 证明: 即使 $A$ 闭, $B$ 也\emph{不必}闭.
\end{enumerate}
\end{exercise}

\begin{exercise}
\label{exer:2.25}
设 $E$, $F$ 与 $G$ 为三个 Banach 空间.

\begin{enumerate}[label=\arabic*., leftmargin=2.2em, itemsep=0.25em]
	\item 设 $T\in\LL(E,F)$, $S\in\LL(F,G)$.证明
	      \begin{equation*}
	      (S\circ T)^\star=T^\star\circ S^\star.
	      \end{equation*}
	\item 假设 $T\in\LL(E,F)$ 是双射.证明 $T^\star$ 是双射, 且
	      $(T^\star)^{-1}=(T^{-1})^\star$.
\end{enumerate}
\end{exercise}

\begin{exercise}
\label{exer:2.26}
设 $E$ 与 $F$ 为两个 Banach 空间, $T\in\LL(E,F)$.设
$\psi:F\to(-\infty,+\infty]$ 为凸函数.假设在 $R(T)$ 中存在某个元素, 使 $\psi$
在该处有限且连续.置
\begin{equation*}
\varphi(x)=\psi(Tx),\qquad x\in E.
\end{equation*}
证明: 对每个 $f\in F^\star$,
\begin{equation*}
\varphi^\star(T^\star f)
=\inf_{g\in N(T^\star)}\psi^\star(f-g)
=\min_{g\in N(T^\star)}\psi^\star(f-g).
\end{equation*}
\end{exercise}

\begin{exercise}
\label{exer:2.27}
设 $E$, $F$ 为两个 Banach 空间, $T\in\LL(E,F)$.假设 $R(T)$ 具有有限余维,
即存在 $F$ 的有限维子空间 $X$, 使得 $X+R(T)=F$ 且 $X\cap R(T)=\{0\}$.

证明 $R(T)$ 闭.
\end{exercise}
